% Generated draft: requires mathematical and source review.
\documentclass[11pt]{article}
\usepackage[margin=25mm]{geometry}
\usepackage{fontspec}
\setmainfont{Latin Modern Roman}
\usepackage{amsmath,amssymb,mathtools,unicode-math}
\setmathfont{Latin Modern Math}
\usepackage{enumitem,needspace,xurl,hyperref}
\setlength{\parindent}{0pt}
\setlength{\parskip}{5pt}
\setlength{\emergencystretch}{3em}
\newcommand{\N}{\mathbb{N}}
\newcommand{\Z}{\mathbb{Z}}
\newcommand{\Q}{\mathbb{Q}}
\newcommand{\R}{\mathbb{R}}
\newcommand{\C}{\mathbb{C}}
\newcommand{\F}{\mathbb{F}}
\newcommand{\A}{\mathbb{A}}
\newcommand{\PP}{\mathbb{P}}
\newcommand{\E}{\mathbb{E}}
\newcommand{\eps}{\varepsilon}
\newcommand{\dd}{\,\mathrm{d}}
\begin{document}
\section*{UnsolvedMath ID: 30001311}
\textbf{Source problem code:} catalogue entry
\subsection{Definitions and mathematical statement}
Let \((B,d_B)\) be a spherical building with its standard spherical metric, normalized so that apartments are unit spherical Coxeter complexes. For a subset \(C\subseteq B\), use the restricted metric \(d_C=d_B|_{C\times C}\). Interpret convexity as \(\pi\)-convexity: every geodesic segment in \(B\) of length strictly less than \(\pi\) whose endpoints lie in \(C\) is contained in \(C\). Closedness is with respect to \(d_B\). Here a subbuilding means a subset carrying a spherical building structure inherited from \(B\); the supplied record does not specify the precise convention for inherited apartments or building type.

For nonempty \(C\), define its intrinsic circumradius by
\[
\operatorname{rad}_C(C):=\inf_{x\in C}\sup_{y\in C}d_C(x,y).
\]
Let
\[
\operatorname{Isom}(C):=\{g:C\to C\mid g\text{ is bijective and }d_C(gx,gy)=d_C(x,y)\text{ for all }x,y\in C\}.
\]
A fixed point of this group means a common fixed point in \(C\).

Under these conventions, the supplied question is whether, for every spherical building \(B\) and every nonempty closed convex subset \(C\subseteq B\) that is not a subbuilding, both assertions hold:
\[
\operatorname{rad}_C(C)\leq\frac{\pi}{2},
\qquad
\exists c\in C\ \forall g\in\operatorname{Isom}(C),\quad g(c)=c.
\]
The record does not explicitly assume nonemptiness or define \(\operatorname{rad}(C)\). Thus this is a conditional precise formulation using the nonempty, intrinsic-radius convention. If centers are instead allowed in \(B\), the relevant quantity is
\[
\operatorname{rad}_B(C):=\inf_{x\in B}\sup_{y\in C}d_B(x,y),
\]
and the intended radius assertion must be confirmed. The record supplies no thickness, dimension, rank, or building-type restriction, and does not require \(C\) to be a simplicial subcomplex. These unresolved conventions must be checked before treating this draft as a definitive formulation.
\subsection{Short English statement}
For a nonempty closed convex subset of a spherical building that is not itself a subbuilding, ask whether its intrinsic circumradius is at most \(\pi/2\) and whether all its isometries fix a common point in the subset. The radius convention, nonemptiness, and precise subbuilding convention require confirmation. The supplied record labels the problem partially solved; this draft does not independently verify that assessment.
\subsection{Sources}
\begin{itemize}
\item Local UnsolvedMath registry: \texttt{data/unsolvedmath/problems.json}, with \texttt{catalog/problems.json} as the additional-record registry; supplied permanent ID \(30001311\), exact source code \texttt{catalogue entry}, title \emph{Radius and Fixed Points of Convex Building Subsets}. Collection: Oberwolfach Reports Open Problems. The record gives proposed year 2009 and names no individual proposer.
\item Supplied report citation: \emph{Differentialgeometrie im Großen (2010)}, Oberwolfach Reports, DOI 10.4171/owr/2009/35, pp.\ 1943--1972. \url{https://doi.org/10.4171/owr/2009/35}. The background attributes the problem to workshop open-problem discussions documented by Oberwolfach Reports.
\item Supplied literature reference, labeled ``Ramos-Cuevas theorem'': \url{https://arxiv.org/abs/1408.3017}.
\item Supplied literature reference, labeled ``Kleiner--Leeb rigidity questions'': \url{https://math.nyu.edu/~bkleiner/convex.pdf}.
\item Supplied literature assessment checked at \texttt{2026-08-21}: status ``partially solved.'' The assessment states that the radius and fixed-point assertions are proved in several cases, including low dimensions, type A, type D, and the thick-building center conjecture, while the full general formulation remains unsettled. This is supplied dated curation derived from the OWR record, not an independently verified literature assessment; its time-sensitive status and cited claims require independent checking.
\item The supplied statement-status label is \texttt{corrected\_verified}. Its verification note says that the mathematical content was checked against the official Oberwolfach report and that notation, context, and question form were normalized without changing the source-level claim. This is a supplied verification claim, not verification performed for this draft.
\end{itemize}
\end{document}
